\chapter{Metodología y Ciclo de Vida}
\section{Ciclo de Vida}

Para la realización del TFG se va a utilizar un ciclo de vida \textbf{Incremental}. Se va a utilizar este enfoque para abordar las tareas de forma modular a la hora de integrar los servicios externos. 


\subsection{Incrementos}
El proyecto se va a dividir en los siguientes incrementos: 

\subsubsection{Incremento I Definición de requisitos y de tecnologías}

En el primer incremento se van a definir los requisitos principales que rigen la aplicación y las tecnologías que se van a usar en el resto de incrementos.

\subsubsection{Incremento II Backend y Frontend de la aplicación}

En este incremento se va a trabajar con la parte central del proyecto, el core, donde se va a encontrar el ciclo de vida del usuario. Esta parte incluye la base de datos relacional de los usuarios y el periodo de pertenencia en el grupo.

Las funcionalidades que se buscan implementar en este ciclo son la gestión de roles y el dashboard visual para el administrador. 

\notaseguridad{Hay que destacar que en el sistema la base de datos que se construye va a ser el \textbf{Single Source of True} \cite{atlassian_ssot},esto se debe a que va a ser el punto centralizado donde está contenido toda la información y todos los permisos que tiene un usuario. De esta manera se puede minimizar el vector de ataque con cuentas de usuarios que ya no están o servicios que no deberían seguir funcionando.}


\subsubsection{Incremento III Infraestructura}

En este incremento se van a gestionar los conflictos que se encuentran en los servidores con los usuarios nativos. Las funcionalidades principales a implementar son los conectores con los servidores y la gestión de recursos de cada usuario. Además se mantendrá un registro de autorizaciones de los servicios que tiene cada una de las personas. 

Los servicios al personal tendrán cinco estados: 

\begin{itemize}
  \item \textbf{Activo}: el servicio se encuentra levantado y operativo.
  \item \textbf{Desactivado}: el servicio se encuentra detenido y los datos se mantienen.
  \item \textbf{Eliminado}: el servicio ha sido eliminado y los datos han sido borrados.
  \item \textbf{Producción}: El servicio se encuentra en producción.
  \item \textbf{Error}: El servicio se encuentra en estado de error.
\end{itemize}


\subsubsection{Incremento IV Integración con servicios externos}

En este incremento se busca trabajar con APIs de terceros para poder gestionar servicios asociados. Dentro de este paso estarían las conexiones con el GitLab, Microsoft Teams y los servicios de VPN.


\subsubsection{Incremento V Auditoría y Cumplimiento}
En este incremento se busca la implementación de un módulo de trazabilidad para medir la robustez. Debido a la estimación en la memoria intermedia de este incremento se debe rescindir. 

Para ello, se tomarán como referencia las directrices del \textbf{NIST Cybersecurity Framework (CSF)}, específicamente en las funciones de \textit{Protect} (Protección) y \textit{Detect} (Detección). Se implementará un sistema de \textit{logging} exhaustivo que garantice la trazabilidad de todas las operaciones de provisión de servicios, cumpliendo con las recomendaciones del estándar \textbf{NIST SP} sobre gestión de registros de log. \cite{nist_csf}

Debido a los incrementos previos y al número de tareas, en la planificación de la entrega intermedia se elimina de producción este incremento.


\subsection{Estructura de los incrementos}
El primer incremento y el cuarto incremento, son los único exentos de esta estructura. El primero está exento ya que se definen globalmente los requisitos que se deben cumplir en el sistema y las tecnologías que se van a usar en primera instancia.  El cuarto se encuentra exento ya que se explican posibles soluciones con marco teórico y la integración de gitlab tienen su propia información. 

\subsubsection{Especificación y modelado} 

En esta etapa se revisan los requisitos y las tecnologías que se van a usar en busca de detalles más específicos. Además, se van a diseñar los diferentes casos de uso del incremento. 

\subsubsection{Diseño de la solución}

Diseño de las pantallas que se van a utilizar y los botones que controlan las diferentes pantallas si aplica. Definición de la lógica que se va a utilizar en la implementación. 

\subsubsection{Implementación y codificación}
Programación de la lógica de negocio de cada uno de los incrementos. 

\subsubsection{Validación y cierre}
Pruebas de cada una de las funcionalidades individuales y pruebas de integración con los ciclos anteriores. Documentación técnica para la memoria del TFG.

% En el apartado de validación y cierre se van a documentar casos de ataque asociados al incremento para verificar la seguridad de la aplicación o servicio.
